\documentclass[11pt]{article}
\usepackage[a4paper,margin=25mm]{geometry}
\usepackage{amsmath,amssymb,amsthm,hyperref}
\newtheorem{theorem}{Theorem}
\newtheorem{lemma}{Lemma}
\newcommand{\Rea}{\operatorname{Re}}
\title{Hurwitz stability of the type A minuscule Tur\'an polynomials}
\author{Principia Math In Progress}
\date{22 September 2026}
\begin{document}
\maketitle

\noindent\textbf{Task 1 (MathDB 359743).}
For $n\ge1$ let
\[
 N_n(x)=\frac1{4n+2}\sum_{k=1}^{n-1}
 k(n-k)\binom{2n+2}{2k+1}x^k.
\]
Ding and Zeng conjecture that $N_{n+1}^2-N_{n+2}N_n$ is weakly
Hurwitz stable for every $n\ge1$ \cite[Conjecture 15]{DZ}.
We prove this and exclude every nonzero zero from the imaginary axis.

\begin{theorem}
For every integer $m\ge2$, put $D_m=N_m^2-N_{m+1}N_{m-1}$.
The zero at $x=0$ has multiplicity exactly two, and every other zero
has strictly negative real part. In particular, the conjecture holds.
\end{theorem}

Let $T_d,U_d$ denote the classical Chebyshev polynomials, with
$U_0(w)=1$, $U_1(w)=2w$, and
$U_{d+1}(w)=2wU_d(w)-U_{d-1}(w)$.

\begin{lemma}\label{lem:bound}
If $d\ge1$ and $|w|>1$, then $|U_d(w)|>d+1$.
If $d\ge2$, $|w|=1$ and $w\ne\pm1$, the same strict inequality holds.
\end{lemma}
\begin{proof}
The recurrence gives
$U_d(\cos\theta)=\sin((d+1)\theta)/\sin\theta$ and $U_d(1)=d+1$.
Its $d$ simple roots are $\cos(j\pi/(d+1))$, $1\le j\le d$.
Pairing the nonzero roots as $\pm r_j$, where $0<r_j<1$, gives
\[
 U_d(w)=2^d w^\epsilon\prod_{j=1}^{\lfloor d/2\rfloor}(w^2-r_j^2),
 \qquad \epsilon=d\bmod2.
\]
For $\rho=|w|\ge1$, the reverse triangle inequality yields
\[
 |U_d(w)|\ge 2^d\rho^\epsilon
 \prod_j(\rho^2-r_j^2)=U_d(\rho).
\]
This product is strictly larger than $U_d(1)=d+1$ if $\rho>1$.
On $\rho=1$, equality in a factor requires $w^2$ to be positive real.
For $d\ge2$ there is at least one such factor, so equality is impossible
unless $w=\pm1$.
\end{proof}

\begin{proof}[Proof of the theorem]
We use the explicit identity from the proof of
\cite[Theorem 14, p.~10]{DZ}, with its index shifted:
\begin{align}
64D_m(z^2)={}&(1+z)^{4m}+(1-z)^{4m} \notag\\
 &-2\bigl(1+(8m^2-2)z^2+z^4\bigr)(1-z^2)^{2m-2}.
\label{eq:source}
\end{align}
This formula is an attributed input. The new step is to recognize its
Chebyshev factorization.
For $x\ne1$ put $w=(1+x)/(1-x)$ and choose $z$ with $z^2=x$.
Writing $q=(1+z)/(1-z)$, one has $(q+q^{-1})/2=w$, hence
\[
 (1+z)^{4m}+(1-z)^{4m}
 =2(1-x)^{2m}T_{2m}(w).
\]
Moreover,
\[
 \frac{1+(8m^2-2)x+x^2}{(1-x)^2}
 =1+2m^2(w^2-1).
\]
The classical identities $T_{2m}=2T_m^2-1$ and
$T_m^2-1=(w^2-1)U_{m-1}^2$ therefore turn
\eqref{eq:source} into
\begin{equation}\label{eq:key}
 D_m(x)=\frac{x(1-x)^{2m-2}}4
 \left[U_{m-1}\!\left(\frac{1+x}{1-x}\right)^2-m^2\right].
\end{equation}
All identities here are polynomial identities after clearing denominators;
the square-root choice has no effect.

If $\Rea x>0$ and $x\ne1$, then
\[
 |w|^2-1=\frac{4\Rea x}{|1-x|^2}>0.
\]
Lemma~\ref{lem:bound} implies $|U_{m-1}(w)|>m$, so
the bracket in \eqref{eq:key} cannot vanish. Its prefactor is nonzero.
At the excluded point $x=1$, substituting $z=1$ in
\eqref{eq:source} gives $D_m(1)=2^{4m-6}>0$.
This proves weak Hurwitz stability.

For a nonzero purely imaginary $x$, $|w|=1$ and $w\ne\pm1$.
When $m\ge3$ the second part of the lemma again excludes a zero.
For $m=2$, direct evaluation gives $D_2(x)=4x^2$.
Finally $[x]N_m=m(m^2-1)/3$ (also valid for $m=1$), and direct
algebra yields
\[
 [x^2]D_m=
 \left(\frac{m(m^2-1)}3\right)^2
 -\frac{(m+1)((m+1)^2-1)(m-1)((m-1)^2-1)}9
 =\frac{m^2(m^2-1)}3>0.
\]
Each $N_m$ has zero constant term. Thus $x^2$ divides $D_m$
and this coefficient proves that its multiplicity is exactly two.
\end{proof}

\noindent\textbf{Verification and scope.}
The proof covers every positive integer in the source conjecture, including
the initial case $N_2^2-N_3N_1=4x^2$.
The accompanying exact-arithmetic verifier checks the factorization from
the defining binomial coefficients, the two symbolic modulus identities,
and finite Hurwitz tests. These finite checks support the written infinite
proof. The source identity and classical Chebyshev facts are credited;
no independent peer review or formal verification is claimed.

\medskip
\noindent\textbf{Proposed Tasks 2--5 (not additional solutions).}
\begin{enumerate}
\setcounter{enumi}{1}
\item Determine the asymptotic size of the largest real part among the
nonzero zeros of $D_m$ as $m\to\infty$.
\item Determine the weak limit of the normalized zero-counting measures
of $D_m(x)/x^2$ after an appropriate scaling.
\item Classify the real parameters $c$ for which
$N_m(x)^2-cN_{m+1}(x)N_{m-1}(x)$ is weakly Hurwitz stable for every $m\ge2$.
\item Determine for which integers $r\ge2$ all polynomials
$N_m(x)^2-N_{m-r}(x)N_{m+r}(x)$, $m\ge r+1$, are weakly Hurwitz stable.
\end{enumerate}

\begin{thebibliography}{1}
\bibitem{DZ} M.-J. Ding and J. Zeng,
\emph{Some New Results on the Minuscule Polynomials of Type A},
Journal of Integer Sequences \textbf{28} (2025), Article 25.1.2.
\href{https://cs.uwaterloo.ca/journals/JIS/VOL28/Zeng/zeng26.pdf}{Published version}.
Preprint: \href{https://arxiv.org/abs/2308.16782}{arXiv:2308.16782v3},
\emph{Real-rootedness of the type A minuscule polynomials}, 6 July 2024;
the target there is Conjecture 4.2.
\end{thebibliography}
\end{document}
